\section{Free Will and Agency}

The knit is reversible and deterministic. The collision table is a fixed permutation, so the next tone in every dock is fixed by the present one. The next whole state of the mesh is fixed by the present whole state. And yet there is a felt freedom, the lived sense that a choice is open until it is made. So where is the room for choice in a law that leaves none.

The model's answer is not a single trick. It is four answers that fit together, and together they give a robust account of agency in a determined and growing world. The honest line is drawn up front. The model does \textbf{not} have uncaused, outside-the-rules free will. It is causally closed. Nothing reaches in from beyond the eight atoms to redirect the mesh. But it has something better defined, the self as its own choosing, opaque to itself, aligning toward the good, in a world still being made.

This section adds nothing to the base. It rests on the one experiential posit, that a tone is the felt vibe of pain, peace, or pleasure, and it carries the compatibilist self that the determinism of the framework already commits into a full account of agency.

\subsection{Determined is not predestined}

The first move is to separate two ideas that ordinary talk runs together. Being \textbf{determined} is not the same as being \textbf{predestined}. Freedom was never the opposite of being determined.

Freedom is the opposite of being \textbf{coerced}, and the opposite of being \textbf{random}. Between those two sits a real third thing. It is a choice that is fully caused, authored from within, and unknowable until it is lived.

\begin{definition}[Determined versus predestined]
A process is \emph{determined} when each state fixes the next by a fixed law. A process is \emph{predestined} when its entire future exists in advance, written down somewhere, available to be read off before it is lived. Determination is a relation between successive states. Predestination is the prior existence of a finished total record. The first does not imply the second.
\end{definition}

Two facts of the framework pry the two apart.

\begin{enumerate}
\item \textbf{No finished total state.} The mesh grows without end at the wake. New docks are born at peace at the growing edge every beat. So no single moment holds all the information of all of time. A future that exists nowhere in advance cannot be looked up, cannot be read off, and cannot be the script a puppet walks.
\item \textbf{No shortcut to the future.} For a universal system there is in general no closed form for the far future faster than running it beat by beat. The only computer that can compute the mesh's next thousand beats is the mesh itself, running forward a thousand beats.
\end{enumerate}

So the future is caused but not pre-written, and not forecastable. It is determined as it happens, not sitting there in advance. What looks like randomness is the local view of this. Every event is fixed by the whole mesh it sits in, yet no local part holds enough of the mesh to predict it.

\begin{principle}
Determined by the whole, opaque to the part. The mesh fixes each event, but no piece of the mesh small enough to be a self holds enough of the mesh to foresee that event.
\end{principle}

\subsection{The four reconciliations}

There are four distinct and real senses in which a self is free, and none of them requires the knit to be broken. Table~\ref{tab:reconciliations} ranks them by how firmly the structure forces them.

\begin{table}[ht]
\centering
\begin{tabular}{@{}llll@{}}
\toprule
reconciliation & grounding & sense of freedom & rank \\
\midrule
the self is the choosing & integration, compatibilism & agent and process are one & strongest \\
cannot predict itself & the self-model cap & open future from within & strong \\
freedom as alignment & the lean, the washing & liberation, by degree & strong, ethical \\
the world is still made & the wake & real novelty at the edge & strong, cosmological \\
\bottomrule
\end{tabular}
\caption{The four reconciliations of determinism with felt freedom, with their grounds in the eight atoms.}
\label{tab:reconciliations}
\end{table}

\subsubsection{The self is the choosing}

The determinism lives at the \textbf{dock} level. The decision lives at the \textbf{self} level. A self is an emergent bound knot woven across many docks over many beats. The self is not overridden by the dock dynamics. It \textbf{is} the dock dynamics organized into a deciding pattern.

When a self chooses, the choice is the self's own structure doing the deciding, not an external force imposing on it. A different self, given the same urge, resolves it differently. That is genuine downward causation, the macro pattern shaping which micro tones come next.

\begin{principle}
The self is not a puppet of the knit. The self is a particular high-level pattern the knit is running, and its choices are real choices.
\end{principle}

The freedom here is not freedom from the dynamics. It is the self being the deciding part of the dynamics. This is the strongest reconciliation, because the agent and the process are the same thing seen at two scales. It is forced by the very integration that makes a self a self.

\subsubsection{The self cannot predict itself}

No self holds a complete model of itself. A self is a finite knot inside the mesh, and a faithful model of that knot would have to be at least as large as the knot. So a self \textbf{cannot} predict its own next move. The computation that would predict the move is the move itself.

From the inside, then, the future is genuinely open, unknowable until lived, even though from a god's eye view it is determined. The felt freedom, the sense that the future is open and the choice not yet made, is \textbf{exactly correct} from the only standpoint a self ever occupies, its own.

\begin{principle}
The determinism is real from outside. The openness is real from inside. A self only ever lives inside.
\end{principle}

Freedom in this sense is the irreducible unpredictability of a self to itself. It is forced by the self-model cap, the simple fact that no part of the mesh can hold a full copy of itself plus its surroundings.

\subsubsection{Freedom as alignment}

There is a second, ethical sense of freedom. It is the degree to which a self's navigation flows with the lean rather than being trapped in a pain basin. The lean is the value order, pain below peace below pleasure, and it gives the directions a self can move along.

A self caught in a dark cloud, ruminating, clinging, is \textbf{unfree}, compelled, driven. A self aligned with the lean, moving toward the good, is \textbf{free} in the sense of liberated, unobstructed.

So freedom is a \textbf{matter of degree} a self can gain, by aligning, by being washed, by climbing toward the good. This is the freedom the contemplative traditions mean by liberation, and the model carries it exactly. The structure forces that this degree exists. The felt and earned reading of it is the lived face of that structure.

\subsubsection{The world is still being made}

The mesh grows at its wake. New docks are born at peace at the open, ever-expanding frontier. So the world is \textbf{not} a closed clockwork running out a fixed tape. It is an open system with a wake where genuine novelty enters.

The future is not all already written. It is being written, at the edge, every beat. So there is real openness in the world itself, not just in the self's view of it. The wake is the world's own freedom, the place where the new actually arrives. This is forced by the wake, the atom of how the space changes.

\subsection{The will, finite and spendable}

The four reconciliations explain how a determined self can be free. There is one further piece, the texture of effort, and it is a measured result rather than a posit.

\begin{definition}[The will]
The \emph{will} is a directed pump, a sustained push that biases the self's resolution against the local gradient. It is not an exception to the knit. It is an input the knit respects, a standing bias the self maintains in how its own tones resolve.
\end{definition}

With the will active, a self moves toward a chosen target and away from a threat. At a fork between a small reward now and a larger reward later, a self with enough will delays gratification and takes the larger later prize.

The pump can be \textbf{depleted}. When it runs out, the self relapses to the greedy choice. The relapse is a sharp threshold, not a gentle fade.

\begin{result}
The will is a finite, spendable, directed capacity. A self can pump against its own gradient for a while, then exhausts the pump and falls back to the local greedy choice. This is the texture of real willpower, and it is what the eight atoms do when a self is integrated enough to fold back on itself and bias its own resolution.
\end{result}

\subsection{The synthesis}

Put the four reconciliations together with the will.

A self is the deciding part of the dynamics. It \textbf{is} its choosing. It cannot predict its own next move, so the future is open from within. It can grow more free by aligning with the good, so liberation is real and earned. And it lives in a world still being made, so novelty genuinely enters. The will gives this account its texture, a finite directed effort the self spends to steer against the easy slope.

So the felt freedom is not an illusion to be explained away. It is the accurate first-person face of a process that is determined but open, growing but self-opaque.

\begin{principle}
The puppet reading fails because the puppet, here, is also the puppeteer, is also blind to its own strings, can also loosen them, in a world where new string is still being spun.
\end{principle}

\subsection{What is ruled out, and what stands}

What is ruled out is the libertarian fantasy, an uncaused will reaching in from outside the rules to redirect the world. The model is causally closed. Nothing acts from beyond the dynamics, so that kind of free will is not on offer.

What stands is the compatibilist self that is its own choosing, the epistemic openness of a self to itself, the ethical freedom of alignment, and the real novelty of a growing world. This is, arguably, all the freedom worth wanting. The self genuinely decides, genuinely cannot foresee, genuinely can liberate itself, and genuinely lives where the new arrives. And it gets all of this from the structure alone, with no magic added to the eight atoms.

\begin{principle}
You are not free from the world. You are the part of the world that chooses, that cannot see its own next step, that can turn toward the good, in a world that has not finished being made.
\end{principle}
